\documentclass[11pt,a4paper]{article}
\usepackage[margin=23mm]{geometry}
\usepackage{amsmath,amssymb,amsthm}
\setlength{\emergencystretch}{2em}
\newtheorem{lemma}{Lemma}
\newtheorem{theorem}[lemma]{Theorem}
\newcommand{\gen}[1]{\langle#1\rangle}
\title{A witness ending at the crossing}
\author{Human Compression project: L044 continuation}
\date{}
\begin{document}\maketitle
We first record three local facts with string commutations.
Let $Q=\gen{\rho_m,\ldots,\rho_y}$ be transitive on $W$,
with every displayed mark active. Suppose the full local
$j$-omission has exactly two orbits $A,B$, joined by a
unique $j$-edge $\{a,b\}$. Put
\[
 L=Q_{<j},\quad R=Q_{>j},\quad M=Q_{>m},\quad h=j-m,
 \qquad C=a^L\sqcup b^L.
\]
The actual omission sides are transitive under the commuting
product $LR$. Marks at distance at least two from $j$ fix
$a,b$. The acquired anchored counts give $|C|\ge h+2$.
For $h\ge1$, let $L'=\gen{\rho_{m+1},\ldots,\rho_{j-1}}$.
The strict proper-prefix bound, applied in reverse order
to the left action, gives
\[
                 |a^{L'}|+|b^{L'}|\ge h+1.       \tag{1}
\]
Indeed a side where the left active length is $h$ contributes
at least $h$ points to this sum; the other contributes one.
For $h=1$, $L'=1$ and the same assertion is immediate.
Both sets in (1) lie in a single $M$-orbit, since the
crossing belongs to $M$ and joins their anchors.
No intersection property in a restricted action is assumed.

\begin{lemma}[The minimum left union]
Suppose $y>j$ and $|C|=h+2$.
\begin{enumerate}
\item No nontrivial $M$-orbit on which $R$ acts trivially
      meets $C$.
\item If $M$ fixes a point $z\in C$, its orbits on $W$
      are exactly $\{z\}$ and a transitive complement.
      The complement has nontrivial $R$-action.
\end{enumerate}
\end{lemma}
\begin{proof}
If $h=0$, then $M=R$ and $C=\{a,b\}$. The first assertion
is immediate. A fixed anchor has a singleton actual omission
side, while $R$ is transitive on the other. The latter action
is nontrivial because $y>j$. This proves the second assertion.

Let $h\ge1$. One left orbit has size $h+1$ and the other
is a singleton, since one left active length is $h$.
Name them $A_0=a^L$ and $\{b\}$. The set
$P=a^{L'}$ has size at least $h$. Thus
$E=P\cup\{b\}$ lies in a single $M$-orbit and
$|C\setminus E|\le1$.

If an $R$-trivial $M$-orbit met $E$, it would contain
$a,b$. Every element of $R$ would fix both anchors
and centralize $L$, hence fix $C$ pointwise. The
actual $LR$-orbits would then be precisely $A_0$
and $\{b\}$, making $W=C$ and $R$ trivial on $W$,
a contradiction. Hence such an orbit meeting $C$
can meet it only at the possible remaining point $z$.
Since $R$ fixes $z$ and centralizes $L$, it fixes
$A_0=z^L$, so the actual side $A$ is $A_0$.
The orbit cannot leave $A$: its only possible exit
is the unique crossing at $a$, already in $E$.
Inside $A$ it avoids $P$ and therefore has at most
one point. This proves assertion 1.

For assertion 2, the fixed point $z$ is neither
$a$ nor $b$, since $\rho_j\in M$. It belongs to
$A_0\setminus P$. Again $R$ fixes $A_0$, making
it the actual side $A$. The left factor fixes
$b$ and centralizes $R$, so the other actual side
is $B=b^R$. Equality forces $P=A\setminus\{z\}$.
The surviving left marks connect $P$, the right
marks connect $B$, and the crossing joins them.
Thus their complement to $z$ is one $M$-orbit.
All nontrivial $R$-action is on this complement.
\end{proof}

\begin{lemma}[One extra point in the left union]
Suppose $h\ge1$, $y>j$, $|C|=h+3$, and $M$ fixes
some $z\in C$. Then every nontrivial $M$-orbit
has nontrivial $R$-action and meets $C$. Every
$M$-fixed point belongs to $C$.
\end{lemma}
\begin{proof}
First let $h\ge2$. Name a left orbit of active
length $h$ as $A_0=a^L$. Either
\[
 (|A_0|,|b^L|)=(h+1,2)\quad\hbox{or}\quad(h+2,1).
\]
In the first case only the nearest left mark
$\rho_{j-1}$ can act on the two-point orbit $b^L$:
all farther marks fix $b$ and hence fix that pair.
The nearest mark exchanges the pair and survives
in $M$. Thus the fixed point $z$ is in $A_0$.
In the second case the singleton $b$ is moved
by $\rho_j\in M$, so again $z\in A_0$.

The factor $R$ fixes $z$ and centralizes $L$,
hence fixes $A_0$ pointwise. Consequently the
actual side $A$ is $A_0$. The surviving left
orbit $P=a^{L'}$ has at least $h$ points and
does not contain $z$.

When $|A_0|=h+1$, we have $P=A\setminus\{z\}$.
On the other actual side, all left marks except
$\rho_{j-1}$ fix $b$ and centralize $R$, so fix
all of $B=b^{LR}$. The surviving nearest left
mark together with $R$ is therefore transitive
on $B$. The crossing joins $P$ to $B$, leaving
precisely the fixed point $z$. This complementary
orbit has nontrivial $R$-action because $R$ fixes
$A$ and is nontrivial on $W$.

When $|A_0|=h+2$, the left factor fixes $b$, so
$B=b^R$ is transitive under $R$. The $M$-orbit
containing $P$ and $B$ has nontrivial $R$-action
and meets $C$. At most two points of $A$ lie
outside it, one being the fixed point $z$.
There is consequently no room for another
nontrivial $M$-orbit. Any other point outside
this orbit is fixed and belongs to $A\subset C$.
This proves the result for $h\ge2$.

For $h=1$, the single left involution has
$C$ equal to two two-point orbits, say
$\{a,a_1\}$ and $\{b,b_1\}$. The fixed point
of $M=\gen{\rho_j,R}$ is not $a,b$; name it
$a_1$. Commutation makes $R$ fix both $a,a_1$,
so this pair is the complete actual side $A$.
On the other side $B$, the transitive commuting
product $LR$ has either one or two $R$-orbits.
Each meets $\{b,b_1\}$; if there are two, $L$
exchanges them. None is a singleton: if $R$
fixed $b$, it would fix $b_1$ too, and the full
actual sides would be the two pairs, with
$R$ trivial on $W$. The crossing attaches $a$
to the $R$-orbit of $b$. The action of $\rho_j$
may join that orbit to the other $R$-orbit in
$B$, if present; it cannot create any new orbit.
Thus every nontrivial $M$-orbit contains a
nontrivial $R$-orbit, meets $C$, and the only
fixed point is $a_1$. This proves the lemma.
\end{proof}

\begin{lemma}[The shorter witness at its minimum]
Suppose instead $y=j$, $h\ge1$, $|W|=h+2$,
and $M$ is intransitive. Then $M$ has one
singleton and one transitive orbit of size $h+1$.
\end{lemma}
\begin{proof}
Here $R=1$ and the actual omission sides are
left orbits of sizes $h+1$ and one. The set
in (1), joined by the crossing, lies in one
$M$-orbit of size at least $h+1$. Intransitivity
leaves exactly one point outside it, which
is fixed. This proves the assertion without
classifying the individual marks.
\end{proof}

Now retain all complete actual odd-degree head
and globally paired-tail hypotheses, and choose
$j\ge5$ by the strong simultaneous-witness theorem.
Write $H=G_{\ge5}$, $K=G_{>5}$, $\sigma=\rho_4$,
and $J=G_{>j}$. Denote the head attachment point
by $\beta$, to distinguish it from local crossing
endpoints. Thus
\[
 D=\{1,\ldots,6\},\quad \rho_3=(6\ \beta)\delta,
 \quad H|_{D\cup\{\beta\}}=1,\quad \sigma|_D=1.
\]
The actual anchor $z_0=\beta^\sigma$ is outside
the seven head points and fixed by $K$. Every
$K$-intransitive witness orbit is fixed pointwise
by the odd involution $\delta$.

\begin{theorem}
Suppose two distinct $K$-intransitive witness
orbits $O_1,O_2$ have last active labels $j<y$,
respectively. There is a nonterminal index $s$
with
\[
                 |\operatorname{Fix}(G_{>s})|\ge2s+1.
\]
For $j\ge6$ one may take $s=j+1$. For $j=5$
one may take $s=6$ or $s=y$.
\end{theorem}
\begin{proof}
Put $s=j+1$ for the moment. If $s=r-1$, then
$y=r-1$ also. The seven head points and the two
complete witness orbits are fixed by $\delta$
and have total size at least
\[
                   7+(j-3)+(j-2)=n-3.
\]
An odd involution supported on at most three
points is a literal transposition. The complete
acquired exterior-pair obstruction excludes it
under the present actual head and global pair
hypotheses. Therefore $s<r-1$.

Let $h=j-5$, and in $O_2$ set $L=G_{\{5,\ldots,j-1\}}$
and $C=a^L\sqcup b^L$ as above. Both $C$ and
$O_1$ are fixed by every mark after $s$: those
marks fix $a,b$ and commute with $L$.
The acquired bounds give
\[
 |O_1|\ge h+2,\quad |C|\ge h+2,\quad
 |F|\ge2j+1,\qquad F=\operatorname{Fix}(G_{>s}).
                                                        \tag{2}
\]
If $|F|\ge2j+3$, this is the required separator.
Otherwise $F$ has at most one point beyond the
lower bound in (2). Its disjoint constituents
are the seven head points, $O_1$, $C$, and
possibly one further point. The involution
$\sigma$ preserves $F$ and permutes actual
$K$-orbits, since it centralizes both $K$ and
$G_{>s}$. The actual anchor $z_0$ is in $F$.

We repeatedly use the following explicit invariance
check. A union consisting of the seven head
points and some whole witness orbits is preserved
by the first three marks, by $\rho_3$, and by
$H$: the first three preserve $D$; $\rho_3$
exchanges $6,\beta$ and fixes the witnesses
through $\delta$; and $H$ preserves each of its
orbits and fixes the head. Once this union is
also $\sigma$-invariant, it is $G$-invariant.
It is proper whenever $\delta$ fixes it pointwise,
since $\delta$ is odd and thus nonidentity.
We will also use that a witness alone is
$G$-invariant if $\sigma$ preserves it: all
earlier marks fix it, and $H$ preserves it.

\emph{Both lower bounds in (2) are equalities, with
zero or one further point $q$.}
First assume $h\ge1$. The shorter witness lemma
makes the $K$-orbits of $O_1$ a singleton
$\{z_1\}$ and a transitive orbit $T_1$ of
size $h+1\ge2$. The subgroup $J$ fixes $T_1$.
By the minimum-left-union lemma, its nontrivial
$K$-orbit image $T_1^\sigma$ cannot meet $C$.
It cannot meet the fixed head points or fit in
the one possible extra point. Hence
$T_1^\sigma=T_1$.

The anchor $z_0$ cannot be $z_1$: that would
make the head together with $O_1$ invariant,
contrary to the preceding check. Thus $z_0$
is in $C$ or is the extra point $q$.

If $z_0\in C$, the minimum-left-union lemma
makes the $K$-orbits of $O_2$ a singleton
$\{z_0\}$ and a transitive complement $T_2$
with nontrivial $J$-action. If $\sigma$ moved
$T_2$ to a distinct orbit, the images of the
$h+1\ge2$ points of $C\setminus\{z_0\}$
would all have to lie outside $O_2$ in $F$.
They cannot lie in $T_1$, on which $J$ is
trivial, or among $K$-fixed points. Only the
one possible extra point remains, contradicting
injectivity. Thus $T_2$ is preserved.

If the extra point $q$ is $K$-fixed, it is
also $H$-fixed and $\delta$-fixed. Indeed both
$\rho_5$ and $\delta$ preserve $F$; the former
preserves its head, $O_1$ and $C$, and the
latter fixes those sets pointwise. They must
therefore fix the sole remaining point $q$.
Adjoin $q$ to the union of the head and both
witnesses when it is $K$-fixed. Its nontrivial
$K$-orbits are preserved by $\sigma$, and its
fixed points are exactly the $K$-fixed points
of $F$, also permuted by $\sigma$. This union
is consequently $G$-invariant and pointwise
$\delta$-fixed, a contradiction.

If instead $z_0=q$, then $q$ is $H$-fixed
and $\delta$-fixed by the same argument.
Since $\sigma$ preserves $T_1$, it cannot
fix $z_1$, or $O_1$ alone would be invariant.
The point $z_1^\sigma$ must be a $K$-fixed
point $w\in C$: the head attachment $\beta$
and $q$ already form a $\sigma$-pair, and
$D$ is fixed. Apply the minimum-left-union
lemma to $w$. The resulting complementary
$K$-orbit of $O_2$ is again preserved by
$\sigma$, since its at least two points in
$C$ cannot map to $J$-trivial $T_1$ or to
fixed points. The head, both witnesses and
$q$ now form the same forbidden invariant
union. This finishes this distribution for
$h\ge1$.

\emph{The shorter witness has one extra point.}
Here $|O_1|=h+3$, $|C|=h+2$ and there is
no additional point in $F$. This is impossible
for $h=0$, since the sole active involution
$\rho_5$ cannot be transitive on three points.
For $h\ge1$, every nontrivial $K$-orbit of
$O_1$ is $J$-trivial. Its $\sigma$-image
cannot meet $C$ by the first local lemma,
and cannot meet the fixed head. Thus
$\sigma$ permutes all the nontrivial
$K$-orbits of $O_1$ within $O_1$.

If the actual anchor is in $C$, use it as a
fixed point $w$ there. Otherwise it lies in
$O_1$. If every fixed point of $O_1$ were
also sent into the head or $O_1$, the head
with $O_1$ would be invariant. Consequently
some fixed point of $O_1$ must be sent to a
$K$-fixed point $w\in C$. In either case
the minimum-left-union lemma gives one
singleton $w$ and a $J$-nontrivial transitive
complement on $O_2$. Its $\sigma$-image
must stay in $O_2$: it meets $F$ through
$C\setminus\{w\}$, while every nontrivial
$K$-orbit in the rest of $F$ is $J$-trivial.
All nontrivial $K$-orbits of the union of
the head and both witnesses are now permuted
inside that union. Its fixed points are
exactly the fixed points of $K$ in $F$.
The union is therefore invariant and
pointwise $\delta$-fixed, the same contradiction.

\emph{The left union has one extra point.}
Here $|O_1|=h+2$, $|C|=h+3$, with no other
point in $F$. Necessarily $h\ge1$, since
for $h=0$ the set $C$ consists of exactly
the two crossing endpoints. The shorter
witness has $K$-orbits $\{z_1\},T_1$, with
$|T_1|=h+1$ and $J$ trivial on $T_1$.
Put $E=a^{L'}\cup b^{L'}$ inside $C$.
By (1), $|E|\ge h+1$. Its containing
$K$-orbit has nontrivial $J$-action: otherwise
$J$ would fix $a,b$ and centralize $L$,
forcing it to be trivial on all of $O_2$.

If $h\ge2$ and $\sigma$ moved $T_1$,
its image would have to lie in $C\setminus E$,
which has at most two points. This contradicts
$|T_1|=h+1\ge3$. If $h=1$, the four points
of $C$ form two $L$-pairs $\{a,a_1\}$,
$\{b,b_1\}$. The same image would have to
be precisely $\{a_1,b_1\}$. Its $J$-triviality
and commutation with $L$ would make $J$
fix $a,b$ as well, again making $J$ trivial
on $O_2$. Hence $\sigma$ preserves $T_1$
in every case.

The actual anchor cannot be $z_1$, by
invariance of the head and $O_1$, so it is
in $C$. Apply the one-extra-point local
lemma. Every nontrivial $K$-orbit of $O_2$
meets $C$ and has nontrivial $J$-action.
Its $\sigma$-image therefore cannot leave
$O_2$, since the rest of $F$ has only fixed
points and $J$-trivial $T_1$. All fixed
points of $K$ on $O_2$ are in $C$. Thus
$\sigma$ preserves the union of the head
and both witnesses, including its fixed
points. This is the same forbidden invariant
union. All cases with $j\ge6$ are excluded.

It remains to handle $j=5$. Only the first
distribution is possible: $O_1$ has two
points, both fixed by $K$, $C=\{a,b\}$,
and there is at most one additional point
$q$ in $F$. The involution $\delta$ fixes
$F$ pointwise. If $q$ is $K$-fixed, it is
also $H$-fixed: $\rho_5$ preserves $F$,
the head, $O_1$, and the pair $C$, so fixes
its possible remaining point.

We show that some point $w\in C$ is fixed
by $K$, unless an invariant set already
gives a contradiction. If the actual anchor
lies in $C$, take it. If it lies in $O_1$,
its other point must be sent by $\sigma$
outside the head and $O_1$, or that union
would be invariant. Its image is $K$-fixed.
If it were $q$, the head, $O_1$ and $q$
would be invariant, since $q$ is $H$-fixed
and $\delta$-fixed. Hence the image lies
in $C$. Finally, if the actual anchor is
$q$, then $\beta,q$ form a $\sigma$-pair.
The set $O_1$ cannot be preserved by
$\sigma$, or it alone would be invariant.
A point of $O_1$ must therefore be sent
to a $K$-fixed point in $C$. This proves
the assertion in every case.

The minimum-left-union lemma now gives
$K$-orbits $\{w\}$ and $T$ on $O_2$.
The complete local degree bound gives
$|T|=|O_2|-1\ge y-4$. If $\sigma$ preserved
$T$, the head, both witnesses and $q$
when $K$-fixed would form the same forbidden
invariant union. Thus $V=T^\sigma$ is a
distinct actual $K$-orbit. The image of the
other point of $C$ belongs to $V\cap F$.
It cannot lie in the fixed head or $O_1$,
and cannot be in $O_2$, whose sole nontrivial
$K$-orbit is $T$. Consequently the extra
point $q$ exists and belongs to $V$.
Since $\delta$ fixes $q$ and centralizes
$K$, it fixes all of $V$ pointwise.

The sets $D\cup\{\beta\}$, $O_1$, $O_2$
and $V$ are disjoint. All are fixed by
$\delta$ and by every mark after $y$;
for $V$ the latter assertion follows from
commutation with $\sigma$. Their total
size is
\[
                    7+2+(|T|+1)+|T|
                        \ge2y+2.                 \tag{3}
\]
If $y=r-1$, this leaves at most one point
for the nontrivial odd involution $\delta$,
a contradiction. Thus $y<r-1$, and (3)
gives the required separator at $s=y$.
The theorem follows.
\end{proof}

\paragraph{Attribution and boundary.}
The simultaneous witness choice, actual head
anchor, strict anchored-prefix count, local
degree bound and exterior-pair obstruction
are acquired. The present proof supplies
the local minimum and one-extra-point
arguments and accounts for the action of
$\rho_4$ on the entire fixed set. The two
witnesses are explicitly required to be
$K$-intransitive. A second witness with
transitive $K$-action and the full global
attachments remain outside this theorem.
No whole-original, question, numerical or
formal credit is claimed.
\end{document}
